\section{Versioned Scientific Correspondence}
\label{app:correspondence}
\begin{table*}[t]
\centering\small
\setlength{\tabcolsep}{4pt}
\renewcommand{\arraystretch}{1.09}
\begin{tabularx}{\textwidth}{@{}>{\raggedright\arraybackslash}p{.17\textwidth}>{\raggedright\arraybackslash}X>{\raggedright\arraybackslash}X@{}}
\toprule
Binding & Study A: theft-method-freeze/5.0.0 & Study B: native-150 compact controller \\
\midrule
Knowledge & Source-extracted propositions, macro graph, refined local evidence rules and mapped routes & Common deterministic compilation of all three public tenancy templates \\
Online gate & Rule-local true/false/unresolved guard; macro edges do not gate requests & Inherited applicability scopes, obligation condition and separate acquisition permission \\
Evidence state & Held document types; true-guard confirmation; union of mapped routes & Presence, capability/route adequacy, multi-document joint adequacy; per-capability route selection \\
State update & Recompute requests/next action from a changed packet; paired edit, not acquired evidence outcome & Recompute from a new assessment; benchmark measures one planning realization \\
Population & 9 development + 27 held-out pairs; nine shared concepts, four contexts; 72 units total & 60 development/11 families + 90 protected/17 families; 150 claims, three tenancy domains \\
Learned inference & Guard-major batched V5; 1-call Direct, 1-call representation baseline, 2-call Evidence-first & Five 2-call singleton pipelines; two additional controls reuse CasePath state \\
Endpoint & Route-aware signed additions/withdrawals, 27-pair micro scores & Complete native artifact and separate actual immediate-checklist metrics \\
Analysis & Historical family bootstrap, exact family swaps and wrong-family diagnostic; broad gate failed & Equal domain/family/case finite contrasts; family-deletion sensitivity, no significance claim \\
Parity & Recorded-answer replay over 72 units and 36 deltas & Separate implementation; no inherited V5 or six-role/live-app parity claim \\
\bottomrule
\end{tabularx}
\caption{A common source-to-state-to-evidence interface does not make the two
implementations, populations or measurements interchangeable. Study B extends
the evidence question rather than retrospectively re-evaluating V5.}
\label{tab:full-correspondence}
\end{table*}


The architectural interface is source specification, case assessment,
evidence planning, state replacement and a justified next action. It does
not assert identical knowledge, algorithms, estimands or empirical
validation across implementations. Study A is bound to
\path{casepath.theft-method-freeze/5.0.0}, interpreter V4 and evidence
overlay V3. Study B is bound to the compact native-150 producer and its
\texttt{obligation\_control\_v1} and \texttt{evidence\_demand\_v1} controller.
The native shared public template union is not a transfer of the theft
source pack, and the theft paired cases are not documented transformations
of the native-150 claim roster.

\paragraph{State and route semantics.}
V5 judges local source guards and projects every mapped route on active
rules, plus mapped confirmation evidence for true unconfirmed guards,
minus held document types. It does not implement capability-specific or
joint adequacy. The native controller distinguishes applicability from
acquisition, unknown evidence from missing evidence, and per-capability
route selection from route union. These are substantive version differences,
not serialization aliases. A benchmark comparison between them is not
claimed here. The generic state update in Equation~\eqref{eq:loop} must be
implemented through each product's own evidence-admission boundary; neither
static evaluation measures whether an external document was actually
acquired or a real claim resolved.

\paragraph{Source and instruction costs.}
V5's preparation extracts propositions, synthesizes and refines process
rules, and maps capabilities to its closed catalogue using models and
retained source-only corrections. Its reported hidden inference costs
exclude that earlier preparation and development. Study B's public
compilation is deterministic and shared by all arms; no learned source
induction cost or claim is substituted for it. The common medium-reasoning,
4,096-completion-token draft/review limit belongs only to Study B. The
V5 records send temperature zero and use guard-major batching and
stage-specific limits; those settings are not transferred to the native
pipeline, whose temperature is omitted and effective default unknown.

\paragraph{Historical recorded costs.}
The 303 retained V5 held-out receipt files report one logged attempt each:
Direct uses 54 calls, 2,900,625 prompt and 45,678 completion tokens at
\$4.8347335; Graph as context uses 54 calls, 5,345,421 prompt and 42,663
completion tokens at \$8.4114187; Evidence-first uses 108 calls, 259,890
prompt and 66,450 completion tokens at \$1.0630999; V5 uses 87 calls,
1,005,211 prompt and 243,891 completion tokens at \$5.4395890. These are
retained provider-reported charges, not fresh account attestations or the
full cost of source preparation. The arms are not compute matched.
Study B accounts for every physical charge, including invalid outputs,
separately from planning reservations. Its two adopted route checks cost
\$0.0006 once, and \$0.591912 of excluded predecessor benchmark charges
remain in full historical outlay. No receipt or cost is counted twice
because it is replayed or adopted.

\section{Signed-Change Scoring and Its Interpretation}
\label{app:paired-metrics}

For pair $i$, the V5 scorer evaluates each permitted signed reference
realization $G\in\mathcal G_i$ against prediction $P_i=\Delta(D_i)$. It
selects the maximum lexicographic tuple
\[
\left(F_1(P_i,G), |P_i\cap G|, \operatorname{Prec}(P_i,G),
\operatorname{Rec}(P_i,G), -|G|, \operatorname{sort}(G)\right).
\]
The last component is the frozen deterministic tie-break. Precision and
recall use zero for their empty denominators; $F_1$ is one when prediction
and reference are both empty. These conventions differ from some native
empty-set conventions. Micro scores pool the selected realization counts
across the 27 pairs; family means instead average pair $F_1$ within each
of the nine branch concepts.

The family bootstrap resamples nine whole families with replacement on
each of 10,000 draws, retaining all three held-out contexts of each selected
family. The exact family-swap test enumerates 512 sign assignments to the
nine family mean differences; three comparator p-values receive Holm
adjustment. Their archived adjusted values are $0.796875$, $0.796875$ and
$0.12890625$. The wrong-family diagnostic uses eight cyclic shifts and
5,000 sampled derangements, with add-one score $(1+N_{\ge})/(1+5008)$.
The reported maximum is the maximum of this evaluated collection, not
exhaustive maximization over all derangements. None of these historical
procedures is substituted for Study B's finite descriptive analysis.

The frozen broad gate requires, among other conditions, $F_1$ margins of
at least $0.10$, precision at least $0.70$, and Holm-adjusted family-swap
$p\le0.05$ against every comparator. These three conditions fail. The
archived structural attribution and wrong-pairing conditions pass, but
passing them does not turn the failed broad claim into statistical
superiority. A source-linked explanation can accompany an incorrect
request, and a family-specific lookup can pass a correspondence diagnostic.

\section{Native Complete-Artifact Measurements}
\label{app:native-metrics}

The native evaluator remains distinct from the signed-change scorer.
Native critical-evidence recall is criticality-weighted coverage of active
evidence obligations through any complete permitted document alternative.
Evidence-obligation completeness additionally requires valid
 decision--fact--capability--document relations. Valid-chain precision
counts native requested items with such a chain divided by all native
requested items; its empty-set value is one and must be read beside
coverage. Wrong-branch and premature requests, document states, node/edge
scores, branch/path probes and source locators retain their exact native
contracts. Locator agreement does not establish semantic entailment;
probe accuracy does not directly measure inferred case-predicate accuracy.

The emitted endpoint counts unique immediate identities without activating
candidate conditions from the reference scenario. It is available only
when the reference has a unique allowed document set for every active
obligation, complete consistent known document states and the public label
mapping. For predicted and reference sets $P,G$,
\[
F_1=\frac{2|P\cap G|}{|P|+|G|},\qquad U=\frac{|P\setminus G|}{|P|}.
\]
Both-empty $F_1$, empty-prediction precision and empty-reference recall
are one; empty-prediction burden is zero. Failed cells use adverse frozen
unit penalties rather than these conventions: zero for higher-is-better
rates, one for lower-is-better rates. Unavailable reference labels do not
become empty sets or silently reduce a denominator. The full planning
object remains retained separately from the native projection.

For split or complete-corpus population $S$, each scheduled case has weight
$\omega_i=1/(3m_{d(i)}n_{f(i)})$, where $m_d$ is its domain's original family
count and $n_f$ its family's scheduled case count. Combined weights are
recomputed over all 28 families. Conventional differences use
adverse-imputed unit scores. Conservative benefit-oriented differences
instead contribute $-1$ whenever either producer endpoint is unobserved,
including two failures; reference availability is a separate prerequisite.
Report all families and every single-family deletion with remaining weights
renormalized within the affected domain. These ranges are composition
sensitivity, not confidence intervals. They neither replace the full
population nor establish unseen-family generalization.
